% --- BIBLIOGRAFIA ---
\apartat{Bibliografia i Webgrafia}
\markboth{Bibliografia i Webgrafia}{}
\addcontentsline{toc}{section}{Bibliografia i Webgrafia}

\begin{description}[style=unboxed, leftmargin=0cm, itemsep=1.2em]
    \sloppy
    \raggedright

    \item[ARASAAC.] (s.d.). \textit{¿Qué es la comunicación aumentativa y alternativa?}. Centro Aragonés para la Comunicación Aumentativa y Alternativa. \url{https://arasaac.org/aac/es}

    \item[American Speech-Language-Hearing Association (ASHA).] (2021). \textit{Augmentative and Alternative Communication (AAC)}. \url{https://www.asha.org/public/speech/disorders/aac/}

    \item[Beukelman, D. R., \& Light, J. C.] (2020). \textit{Augmentative and Alternative Communication: Supporting Children and Adults with Complex Communication Needs} (5th ed.). Paul H. Brookes Publishing Co.

    \item[Institut Obert de Catalunya.] (s.d.). \textit{Atenció a la comunicació}. Materials del cicle formatiu d'Atenció a Persones en Situació de Dependència. \url{https://ioc.xtec.cat/materials/FP/Recursos/fp_apd_m09_/web/fp_apd_m09_htmlindex/WebContent/u1/a1/continguts.html}

    \item[Mayo Clinic.] (2023). \textit{Autism spectrum disorder - Symptoms and causes}. \url{https://www.mayoclinic.org/diseases-conditions/autism-spectrum-disorder/symptoms-causes/syc-20352928}

    \item[Pagliano, P.] (2012). \textit{The Multisensory Handbook: A Guide for Children and Adults with Sensory Learning Disabilities}. Routledge.

    \item[Sensory Processing Disorder Foundation.] (2024). \textit{About SPD: Sensory Processing Disorder}. \url{https://www.spdstar.org/basic/about-spd}

    \item[UNESCO.] (2020). \textit{Towards inclusion in education: status, trends and challenges}. \url{https://unesdoc.unesco.org/ark:/48223/pf0000373844}

\end{description}
